% Tabel ulasan pelanggan Nasi Gerilya - 20 ulasan (10 positif + 10 negatif/netral)
% Sumber: K05-GF-20260613-data koding (K05-RP-001 s.d. K05-RP-010; K05-RN-001 s.d. K05-RN-010)

\begingroup
\scriptsize
\begin{xltabular}{\textwidth}{|C{0.6cm}|L{1.8cm}|C{0.8cm}|L{2.2cm}|Y|}
\caption{Daftar Ulasan Pelanggan Nasi Gerilya pada GrabFood}
\label{tab:daftar-ulasan-nasi-gerilya-grabfood} \\
\toprule
\textbf{No.} & \textbf{Nama Pelanggan} & \textbf{Rating} & \textbf{Kategori} & \textbf{Isi Ulasan} \\ \hline
\midrule
\endfirsthead

\continuedtabletitle{5}
\toprule
\textbf{No.} & \textbf{Nama Pelanggan} & \textbf{Rating} & \textbf{Kategori} & \textbf{Isi Ulasan} \\ \hline
\midrule
\endhead

\continuedtablefooter{5}
\endfoot

\bottomrule
\endlastfoot

% --- Ulasan Positif ---
\rowcolor{black!8}\multicolumn{5}{|l|}{\textbf{A. Ulasan Positif}} \\ \hline
1 & Muslim M. & 5 & Positif & Keren ni, konsisten rasa \\ \hline
2 & Raymond & 5 & Positif & Ayam popnya wenak dan \textit{creamy} \\ \hline
3 & tasyaa & 5 & Positif & Rasanya enak, sayur \textit{fresh}, porsi agak kebanyakan \\ \hline
4 & Wie W. & 5 & Positif & Langganan keluarga; rasa enak, porsi besar, bersih, \textit{request} dipenuhi \\ \hline
5 & Wie W. & 5 & Positif & Langgananku setiap Minggu, \textit{Best} \\ \hline
6 & ita & 5 & Positif & Ikan gembung dan lele \textit{fresh}; rasanya enak \\ \hline
7 & Lina & 5 & Positif & Semua menu enak, semua \textit{recommended} \\ \hline
8 & Gusti & 5 & Positif & \textit{Best} karena sayur dan lauk dipisah \\ \hline
9 & Derrick & 5 & Positif & Rasanya enak, porsinya besar \\ \hline
10 & Suba & 5 & Positif & Enak banget menurut keluarga \\ \hline

% --- Ulasan Negatif/Netral ---
\rowcolor{black!8}\multicolumn{5}{|l|}{\textbf{B. Ulasan Negatif/Netral}} \\ \hline
11 & albert p. & 1 & Negatif & Jangek siram tidak dikasih; restoran meminta maaf dan berjanji memperhatikan kelengkapan \\ \hline
12 & Talia & 1 & Negatif & Pesan ayam rendang, dikirim ayam goreng \\ \hline
13 & irawaty & 3 & Netral & Udang goreng hanya 3 biji, terlalu mahal, tidak \textit{worth it} \\ \hline
14 & Hendra & 3 & Netral & Pesan ikan kakap tetapi diberikan ikan nila; restoran meminta maaf \\ \hline
15 & siyaga. & 3 & Netral & Makanan enak tetapi kuah yang diminta/dibayar tidak diberikan; restoran merespons \\ \hline
16 & lily & 2 & Negatif & Nasi agak keras, keripik kentang melempem, \textit{packaging} bagus \\ \hline
17 & Khash & 3 & Netral & \textit{Not quite fresh}; \textit{tough meat}; \textit{okay but does not meet expectation} \\ \hline
18 & Hartini & 1 & Negatif & Sayur nangka ada rada basi \\ \hline
19 & Vivi & 4 & Netral & Staf kurang koordinasi; salah \textit{order} ke \textit{driver}; pelanggan keluar ongkos lagi; rasa biasa saja \\ \hline
20 & Fajar & 4 & Netral & Setelah makan langsung diare; restoran memberi respons \\ \hline

\end{xltabular}
\tablesource[\textwidth]{Observasi GrabFood Nasi Gerilya - Glugur Kota, 13 Juni 2026; diolah peneliti (2026).}

\noindent\footnotesize Catatan: Rating bintang dan isi ulasan tidak selalu searah. Kolom kategori didasarkan pada isi komentar, bukan hanya bintang.
\endgroup
